\chapter{The Legendre Transformation}
\label{ch:leg}

The Legendre transformation is much more than an algebraic procedure for
passing from a Lagrangian to a Hamiltonian. It is a general method for
replacing one set of variables by their conjugate variables while preserving
the underlying physical information. It appears throughout physics: in
analytical mechanics, thermodynamics, geometrical optics, field theory, and
constrained dynamics.

There are three closely related ways of understanding it:
\begin{enumerate}
   \item Geometrically: it replaces points on a graph by the slopes of its
       tangent planes.
   \item Physically: it changes which quantities are regarded as independent
       and which are determined by the constitutive or dynamical relations.
   \item In mechanics: it replaces velocities by canonical momenta and, when
       regular, converts the Lagrangian description into a Hamiltonian
        description with a natural symplectic structure.
\end{enumerate}
 The distinction between the second and third points is particularly important:
 the transformation itself is a mathematical construction, whereas its physical
 interpretation depends on the problem.

 \section{Geometry}
 \label{sec:lag-geom}
 Here I describe Legendre transformation from the physicist point of view. More
 mathematical approach is presented in Appendix \ref{ch:lt}.
 \subsection{1D:points versus tangent lines}
Consider a smooth, strictly convex function \(f(x)\).
At each point \(x\), its tangent line has slope
\[
p=f'(x).
\]
The equation of the tangent line is
\[
y=f(x)+p(X-x)
\]
or equivalently
\[
\boxed{y=pX-g(p)}
\]
where
\[
\boxed{g(p)=px-f(x),\qquad p=f'(x).}
\]
Provided the relation \(p=f'(x)\) can be inverted, \(g\) is a function of \(p\).
This function is the Legendre transform of \(f\).

Thus we have two descriptions of the same family of tangent lines.
In the original description, we specify the contact point \(x\) and calculate
the slope \(p\). In the dual description, we specify the slope \(p\) and
calculate the contact point \(x\).
The crucial differential identity is
\[
dg=d(px-f)
=p\,dx+x\,dp-f'(x)\,dx.
\]
Since \(p=f'(x)\), the \(dx\) terms cancel:
\[
\boxed{dg=x\,dp.}
\]
Consequently,
\[
\boxed{x=\frac{dg}{dp}.}
\]
We therefore have the reciprocal relationships
\[
\boxed{
p=\frac{df}{dx},
\qquad
x=\frac{dg}{dp}.
}
\]
This reciprocity is the essential mathematical content of the Legendre
transformation.

\paragraph{The envelope interpretation}
For convex \(f\), every tangent line lies below the graph:
\[
f(x)\geq px-g(p).
\]
Taking the supremum over all slopes gives
\[
\boxed{f(x)=\sup_p\{px-g(p)\}.}
\]
Similarly,
\[
\boxed{g(p)=\sup_x\{px-f(x)\}.}
\]
The graph of \(f\) is therefore the envelope of its tangent lines.
The Legendre transformation converts a pointwise description of a curve into a
description in terms of its supporting tangent lines.
There is no loss of information for a proper closed convex function:
transforming twice recovers the original function.

\subsection{Several variables and the geometry of dual spaces}
Let \(f(x^1,\ldots,x^n)\) be a function on a vector space \(V\). Define
\[
p_i=\frac{\partial f}{\partial x^i}.
\]
Its Legendre transform is
\[
g(p)=p_ix^i-f(x),
\]
where \(x=x(p)\).
The differential relations become
\[
df=p_i\,dx^i,
\]
\[
dg=x^i\,dp_i.
\]
Hence
\[
\boxed{
p_i=\frac{\partial f}{\partial x^i},
\qquad
x^i=\frac{\partial g}{\partial p_i}.
}
\]
There is an important geometric point here.
The variables \(x^i\) belong to the vector space \(V\), whereas the quantities \(p_i\) belong to its dual space \(V^*\).
The product
\[
p_ix^i=\langle p,x\rangle
\]
is the natural pairing between a vector and a covector.
No Euclidean metric is required to define this pairing.
This will become essential when we discuss mechanics on manifolds.

\paragraph{The Hessian and the inverse response}
Differentiating the conjugacy relations gives
\[
dp_i=\frac{\partial^2 f}{\partial x^i\partial x^j}\,dx^j.
\]
If the Hessian is nonsingular, then
\[
dx^i=\frac{\partial^2 g}{\partial p_i\partial p_j}\,dp_j.
\]
It follows that
\[
\boxed{\operatorname{Hess}(g)
=\left[\operatorname{Hess}(f)\right]^{-1}.}
\]
Thus the Legendre transformation exchanges a response law with its inverse.
This is already a physical idea: one may prescribe a variable and determine its
conjugate response, or prescribe the response and determine the original
variable.

\section{Physical meaning in analytical mechanics}
Consider a mechanical system with generalized coordinates \(q^i\), velocities
\(v^i=\dot q^i\), and a Lagrangian
\[
L(q,v,t).
\]
Its canonical momenta are defined by
\[
\boxed{p_i=\frac{\partial L}{\partial v^i}.}
\]
Assume initially that these relations can be inverted to express the velocities
as
\[
v^i=v^i(q,p,t).
\]
The Hamiltonian is defined by
\[
\boxed{
H(q,p,t)=p_iv^i-L(q,v,t).
}
\]
The Legendre transformation is performed only in the velocity variables; \(q\)
and \(t\) remain fixed.

\paragraph{Why replace velocity by momentum?}
A Lagrangian description specifies the system through configuration and velocity:
\[
(q^i,v^i).
\]
A Hamiltonian description specifies it through configuration and momentum:
\[
(q^i,p_i).
\]
These are not simply two arbitrary choices of coordinates.
Velocity describes how the configuration changes in time. Canonical momentum
measures the sensitivity of the Lagrangian to a change of velocity.

For example, if
\[
L=\frac12 mv^2-V(q),
\]
then
\[
p=mv.
\]
The momentum–velocity relationship is determined by the inertial properties of
the system.
But for a charged particle in an electromagnetic field, in Gaussian units,
\[
L=\frac12 m\mathbf v^2
+\frac ec\mathbf A\cdot\mathbf v-e\phi,
\]
the canonical momentum is
\[
\boxed{\mathbf p=m\mathbf v+\frac ec\mathbf A.}
\]
The mechanical momentum \(m\mathbf v\) and canonical momentum \(\mathbf p\)
differ.
This illustrates something important:
The Legendre transformation does not simply replace \(mv\) by a different
symbol. It introduces the variables conjugate to motion as determined by the
complete Lagrangian.

\paragraph{Why does the Hamiltonian represent energy?}
Let us calculate the differential of \(H\):
\[
\begin{aligned}
dH
&=d(p_iv^i-L)\\
&=v^i\,dp_i+p_i\,dv^i
-\frac{\partial L}{\partial q^i}\,dq^i
-\frac{\partial L}{\partial v^i}\,dv^i
-\frac{\partial L}{\partial t}\,dt.
\end{aligned}
\]
The terms containing \(dv^i\) cancel, leaving
\[
\boxed{
dH=v^i\,dp_i
-\frac{\partial L}{\partial q^i}\,dq^i
-\frac{\partial L}{\partial t}\,dt.
}
\]
Therefore,
\[
\boxed{
\frac{\partial H}{\partial p_i}=v^i,
\qquad
\frac{\partial H}{\partial q^i}
=-\frac{\partial L}{\partial q^i},
\qquad
\frac{\partial H}{\partial t}
=-\frac{\partial L}{\partial t}.
}
\]
The Euler–Lagrange equations give
\[
\frac{dp_i}{dt}=\frac{\partial L}{\partial q^i}.
\]
Consequently,
\[
\boxed{
\dot q^i=\frac{\partial H}{\partial p_i},
\qquad
\dot p_i=-\frac{\partial H}{\partial q^i}.
}
\]
These are Hamilton's equations.
But the appearance of energy has a deeper explanation.
The quantity
\[
E_L=p_iv^i-L
\]
satisfies, along the Euler–Lagrange trajectories,
\[
\frac{dE_L}{dt}=-\frac{\partial L}{\partial t}.
\]
Thus for a time-independent Lagrangian,
\[
\boxed{\frac{dE_L}{dt}=0.}
\]
The Legendre expression is precisely the Noether quantity associated with time
translations.
The Hamiltonian therefore represents energy in the standard autonomous
mechanical setting, although the identification requires care for explicitly
time-dependent, gauge-dependent, and parametrization-invariant formulations.

\paragraph{Example: free particle and relativistic particle}
Nonrelativistic particle
\[
L(v)=\frac12 mv^2
\]
\[
p=mv
\]
\[
H(p)=\frac{p^2}{2m}
\]
Relativistic particle
\[
L(v)=-mc^2\sqrt{1-v^2/c^2}
\]
\[
p=\gamma mv
\]
\[
H(p)=\sqrt{m^2c^4+p^2c^2}
\]
Comparison of momentum–velocity relations
Set \(m=c=1\) for the relativistic example and \(m=1\) for the nonrelativistic
example.
Schematic plot on separate natural scales: in the relativistic case, momentum
diverges as the speed approaches \(c\); in the Newtonian case momentum is
linear in velocity.
In both cases, the Legendre transformation replaces velocity by momentum.
However, the relation between them is completely different.
For the nonrelativistic particle,
\[
v=\frac pm,\qquad
H=\frac{p^2}{2m}.
\]
For the relativistic particle,
\[
v=\frac{pc^2}{\sqrt{m^2c^4+p^2c^2}},
\]
and hence
\[
\frac{dH}{dp}=v.
\]
The same duality holds, independently of the particular dynamical law.
This is one reason the Legendre transformation is so universal.

\subsection{Intrinsic geometric meaning: tangent and cotangent bundles}
The most important geometric interpretation in mechanics concerns the
distinction between the tangent bundle \(TQ\) and the cotangent bundle \(T^*Q\)
of configuration space \(Q\).
For a configuration manifold \(Q\),
\[
q\in Q,\qquad v\in T_qQ,\qquad p\in T_q^*Q.
\]
Velocities are tangent vectors. Momenta are cotangent vectors.
The Legendre map is defined intrinsically by differentiating the Lagrangian
along each tangent-space fiber:
\[
\boxed{
\mathbb FL:TQ\longrightarrow T^*Q.
}
\]
More explicitly,
\[
\boxed{
\mathbb FL(q,v)
=\left(q,\frac{\partial L}{\partial v}(q,v)\right).
}
\]
The second component is a covector: it acts on a velocity variation \(w\in
T_qQ\) according to
\[
p(w)=
\left.\frac{d}{d\epsilon}
L(q,v+\epsilon w)\right|_{\epsilon=0}.
\]
This construction requires no metric and no global linear structure on \(Q\).
Bundle interpretation of the Legendre map
The map keeps the base configuration \(q\) fixed while changing the fiber
variable from \(v\) to \(p\).
For example, if
\[
L(q,v)=\frac12 g_{ij}(q)v^iv^j-V(q),
\]
then
\[
p_i=g_{ij}(q)v^j.
\]
In this particular case, the Legendre map is simply the musical isomorphism
\(v\mapsto v^\flat\) defined by the kinetic-energy metric \(g\)(see Appendix
\ref{sec:math-mus}).
For a general Lagrangian, however, it is a nonlinear map between tangent and
cotangent fibers.

\paragraph{Why is the cotangent bundle special?}
The cotangent bundle possesses a canonical one-form
\[
\boxed{\theta=p_i\,dq^i}
\]
and a canonical symplectic form
\[
\boxed{\omega=-d\theta=dq^i\wedge dp_i.}
\]
These structures exist independently of any metric or particular Hamiltonian.
By contrast, the tangent bundle does not possess a canonical symplectic form
without additional structure.
A regular Lagrangian supplies one through the Legendre map:
\[
\theta_L=(\mathbb FL)^*\theta
=\frac{\partial L}{\partial v^i}\,dq^i,
\]
\[
\boxed{\omega_L=-d\theta_L=(\mathbb FL)^*\omega.}
\]
If the Legendre map is invertible, the two symplectic descriptions are
equivalent.
This provides a deeper interpretation:
The Legendre transformation transports the Lagrangian dynamics from
tangent-bundle geometry to the canonical symplectic geometry of the cotangent
bundle.
Strictly speaking, the Legendre map is not itself a canonical transformation
between two phase spaces already equipped with their standard symplectic
structures. Rather, it establishes a symplectic equivalence after the
Lagrangian structure on \(TQ\) has been constructed.

\paragraph{An even deeper connection: contact geometry}
There is also a geometric structure underlying the Legendre transformation
independently of dynamics.
In one dimension, introduce coordinates \((x,f,p)\) with the contact one-form
\[
\alpha=df-p\,dx.
\]
On the graph of \(f(x)\), where \(p=f'(x)\), one has \(\alpha=0\).
Define new variables
\[
g=px-f.
\]
Then
\[
dg-x\,dp=p\,dx-df=-\alpha.
\]
Thus the transformation preserves the contact distribution, up to an overall
sign.
This means that the Legendre transformation is naturally a contact
transformation on the space of first-order contact elements.
The symplectic structure of Hamiltonian mechanics and the contact structure of
thermodynamics are therefore not accidental parallels: Legendre duality is
naturally compatible with both.

\subsection{What happens when the Legendre transformation is singular?}
The standard Hamiltonian construction assumes that
\[
\det\left(
\frac{\partial^2L}{\partial v^i\partial v^j}
\right)\ne0.
\]
If this Hessian is nonsingular, the Legendre map is locally invertible. Such a
Lagrangian is called regular.
If it is globally invertible onto \(T^*Q\), it is called hyperregular.
But an important class of physical theories has singular Lagrangians.
Homogeneous Lagrangians and vanishing Hamiltonians
Suppose \(L(q,v)\) is homogeneous of degree one in velocities:
\[
L(q,\lambda v)=\lambda L(q,v),\qquad \lambda>0.
\]
Euler's theorem gives
\[
v^i\frac{\partial L}{\partial v^i}=L.
\]
Hence
\[
\boxed{
H_c=p_iv^i-L=0.
}
\]
Moreover, differentiating Euler's identity with respect to velocity gives
\[
\frac{\partial^2L}{\partial v^i\partial v^j}v^j=0.
\]
Therefore the Hessian has a zero mode and the Legendre map cannot be inverted
in the usual way.
This is not a mathematical defect. It is the geometric expression of
reparametrization invariance.
A standard example is the relativistic worldline action
\[
S=-mc\int d\lambda\,
\sqrt{\eta_{\mu\nu}\dot x^\mu\dot x^\nu},
\]
with
\[
\eta_{\mu\nu}=\operatorname{diag}(+,-,-,-).
\]
The canonical momenta satisfy the mass-shell constraint
\[
\boxed{\eta^{\mu\nu}p_\mu p_\nu=m^2c^2.}
\]
The canonical Hamiltonian vanishes, but that does not mean the physical energy
vanishes.
Instead, the Hamiltonian formulation uses constraints, for example
\[
H_T=u(\lambda)(p^2-m^2c^2),
\]
where \(u(\lambda)\) is an arbitrary multiplier.
Choosing physical time as the evolution parameter recovers the familiar
nonvanishing relativistic energy
\[
E=\sqrt{m^2c^4+c^2\mathbf p^2}.
\]
Here we encounter a profound distinction:
\begin{itemize}
   \item The canonical Hamiltonian generates evolution with respect to the
       arbitrary worldline parameter.
   \item The physical energy generates translations in a selected physical time
       coordinate.
\end{itemize}
These are not the same notion.

The theory of constrained Hamiltonian systems, particularly Dirac's theory, can
be understood as the extension of Legendre duality to systems where the
velocity–momentum correspondence is not invertible.

\subsection{ Physical interpretation in thermodynamics}
Thermodynamics provides perhaps the clearest demonstration that a Legendre
transformation is a change of physical description, not merely an algebraic
trick.
Consider the internal energy
\[
U=U(S,V,N).
\]
For a simple compressible system,
\[
dU=T\,dS-P\,dV+\mu\,dN.
\]
Its natural independent variables are \(S,V,N\).
But experimentally, we may not control the entropy \(S\). Instead, we place the
system in contact with a reservoir at fixed temperature \(T\).
We therefore replace entropy by its conjugate variable, temperature:
\[
T=\frac{\partial U}{\partial S}.
\]
Define the Helmholtz free energy:
\[
\boxed{F(T,V,N)=U-TS.}
\]
Its differential is
\[
\begin{aligned}
dF
&=dU-T\,dS-S\,dT\\
&=-S\,dT-P\,dV+\mu\,dN.
\end{aligned}
\]
Thus
\[
\boxed{
S=-\frac{\partial F}{\partial T},
\qquad
P=-\frac{\partial F}{\partial V}.
}
\]
The transformation has changed the natural independent variable from \(S\) to
\(T\).
There is also a physical change in the equilibrium principle. At fixed entropy
and volume, equilibrium is characterized by minimization of the appropriate
internal energy. At fixed reservoir temperature and volume, the relevant
potential is \(F\).
The Legendre transformation makes this change of externally controlled
variables explicit.
Notice the sign convention: since the usual thermodynamic transformation writes
\(F=U-TS\), it is the negative of the conventional convex-conjugate expression
\(TS-U\). Thermodynamics commonly uses this sign convention to retain
energy-like potentials.
Similarly, enthalpy
\[
H_{\mathrm{th}}=U+PV
\]
replaces volume \(V\) by pressure \(P\), because its conjugate in \(dU\) is
\(-P\).
This reveals a general principle:
We perform a Legendre transformation when it is more useful to regard the
conjugate response as independently controlled than the original variable.

\paragraph{Two qualifications are worth emphasizing.}
\begin{enumerate}
   \item If the function is not convex, the global Legendre–Fenchel transform
\[
f^*(p)=\sup_x(px-f(x))
\]
does not generally reconstruct the original nonconvex function. Its double
transform \(f^{**}\) gives the closed convex envelope under the standard
hypotheses.
However, mechanics also uses local Legendre transformations for regular but
nonconvex Lagrangians. These are locally valid even where the global convex
construction is not equivalent. One should distinguish the two notions.
   \item The transformation does not by itself introduce new physical laws. For
       a regular Lagrangian, the Lagrangian and Hamiltonian formulations
        express the same dynamics. What changes is the structure in which those
        laws are formulated.
This distinction matters because the Hamiltonian description provides immediate
        access to Poisson brackets, symplectic reduction, canonical
        transformations, and phase-space methods that are less transparent in
        the original velocity formulation.
\end{enumerate}

The first-order Hamiltonian action is
\[
\boxed{
S[q,p]=\int_{t_1}^{t_2}
\left(p_i\dot q^i-H(q,p,t)\right)dt.
}
\]
Variation with respect to \(p_i\) gives
\[
\dot q^i=\frac{\partial H}{\partial p_i}.
\]
Variation with respect to \(q^i\), with the usual fixed configuration
endpoints, gives
\[
\dot p_i=-\frac{\partial H}{\partial q^i}.
\]
For a hyperregular Lagrangian, eliminating \(p_i\) reproduces
\[
S[q]=\int L(q,\dot q,t)\,dt.
\]
Thus the Legendre transformation establishes an equivalence between two
variational descriptions:
\[
\boxed{
\begin{aligned}
&\text{Configuration-space action: }
\int L(q,\dot q,t)\,dt,\\[3pt]
&\text{Phase-space action: }
\int(p_i\dot q^i-H(q,p,t))\,dt.
\end{aligned}
}
\]
The phase-space action separates the geometric term \(p_i\,dq^i\), which
defines the symplectic structure, from \(H\,dt\), which specifies the dynamics.
This is arguably the most significant physical and geometric meaning of the
Legendre transformation in mechanics: it separates the universal geometry of
phase space from the particular dynamical law governing the system.
It also explains why singular Legendre transformations are so consequential.
When invertibility fails, this clean separation requires constraints, and the
resulting phase space may possess gauge directions rather than a nondegenerate
symplectic structure.

\section{Generating Functions, and Canonical Transformations}
In analytical mechanics, Legendre transformations, generating functions, and
canonical transformations are closely connected. However, they are not
identical, and distinguishing them reveals much of the geometry underlying
Lagrangian and Hamiltonian mechanics.
The central idea is that a Legendre transformation changes the variables used
to describe an object, whereas a canonical transformation changes phase-space
variables while preserving their symplectic relationships. Generating functions
provide a way of constructing and representing canonical transformations.
There is an important subtlety: transformations between different types of
generating functions are themselves partial Legendre transformations.

Let us develop these ideas systematically.
Consider a configuration manifold \(Q\), its tangent bundle \(TQ\), and its
cotangent bundle \(T^*Q\).
\paragraph{ Legendre transformation}
Starting with a Lagrangian
\[
L(q^i,v^i,t),
\]
define
\[
p_i=\frac{\partial L}{\partial v^i}.
\]
When this relation is invertible, the Hamiltonian becomes
\[
\boxed{H(q,p,t)=p_iv^i-L(q,v,t).}
\]
The transformation is
\[
\mathbb FL:TQ\longrightarrow T^*Q.
\]
Thus the velocity variables are replaced by their conjugate momenta, with the
configuration \(q\) unchanged.
Notice that \(TQ\) and \(T^*Q\) are different geometric spaces.

\paragraph{Canonical transformation}
A canonical transformation is a map
\[
\Phi:T^*Q\longrightarrow T^*Q
\]
(or, more generally, between two cotangent bundles)
\[
(q^i,p_i)\longmapsto(Q^i,P_i)
\]
that preserves the symplectic form
\[
\boxed{
\omega=dq^i\wedge dp_i.
}
\]
That is,
\[
\boxed{\Phi^*\omega=\omega.}
\]
In local coordinates,
\[
\boxed{
dQ^i\wedge dP_i=dq^i\wedge dp_i.
}
\]
The old and new variables need not have the same physical interpretation
individually. For example, a canonical transformation may mix positions and
momenta.
Unlike the ordinary Legendre map, a canonical transformation relates two
descriptions that are already Hamiltonian.

\paragraph{Generating function}
A generating function is a scalar function from which a canonical
transformation can be obtained by differentiation.
The simplest version is a function
\[
F_1(q,Q).
\]
It generates a canonical transformation through
\[
\boxed{
p_i=\frac{\partial F_1}{\partial q^i},
\qquad
P_i=-\frac{\partial F_1}{\partial Q^i}.
}
\]
Why do precisely these derivatives appear, and why is there a minus sign?
The answer is geometric.

\subsection{Why generating functions produce canonical transformations}
Introduce the canonical one-form on phase space:
\[
\theta=p_i\,dq^i.
\]
With the convention
\[
\omega=-d\theta=dq^i\wedge dp_i,
\]
the one-form \(\theta\) determines the symplectic structure.
For a canonical transformation, consider the difference
\[
\alpha=p_i\,dq^i-P_i\,dQ^i.
\]
Its exterior derivative is
\[
d\alpha
=dp_i\wedge dq^i-dP_i\wedge dQ^i.
\]
Consequently,
\[
d\alpha=0
\]
is equivalent to
\[
dq^i\wedge dp_i=dQ^i\wedge dP_i.
\]
Thus the transformation is canonical precisely when the one-form \(\alpha\),
restricted to the transformation graph, is closed.
Locally, the Poincaré lemma implies that there exists a function \(F\) such
that
\[
\boxed{
p_i\,dq^i-P_i\,dQ^i=dF.
}
\]
If we can regard \(q\) and \(Q\) as independent coordinates on the graph, then
\(F=F_1(q,Q)\), and
\[
dF_1=
\frac{\partial F_1}{\partial q^i}dq^i+
\frac{\partial F_1}{\partial Q^i}dQ^i.
\]
Comparison gives
\[
\boxed{
p_i=\frac{\partial F_1}{\partial q^i},
\qquad
P_i=-\frac{\partial F_1}{\partial Q^i}.
}
\]
These equations are therefore not merely an algebraic recipe.
They express the fact that the two canonical one-forms differ by an exact
differential.

Canonical transformations preserve the symplectic two-form, but they do not
necessarily preserve the canonical one-form itself.
Instead, locally,
\[
\boxed{\theta-\Phi^*\theta=dF.}
\]
The addition of an exact differential changes the canonical one-form without
changing its exterior derivative.
This simple observation explains why generating functions work.
It also anticipates the freedom to add total derivatives to Lagrangians.

\subsection{Generating functions as partial Legendre transformations}
Here is where the Legendre transformation enters for the second time.
Start from a type-1 generating function,
\[
F_1(q,Q),
\]
satisfying
\[
dF_1=p_i\,dq^i-P_i\,dQ^i.
\]
Suppose we prefer to use \(q\) and \(P\) rather than \(q\) and \(Q\) as the
independent variables.
Define
\[
\boxed{
F_2(q,P)=F_1(q,Q)+Q^iP_i.
}
\]
This is a partial Legendre transformation with respect to the \(Q\) variables,
taking account of the sign convention
\[
P_i=-\frac{\partial F_1}{\partial Q^i}.
\]
Differentiating,
\[
dF_2
=dF_1+Q^i\,dP_i+P_i\,dQ^i.
\]
Substitution yields
\[
dF_2=p_i\,dq^i+Q^i\,dP_i.
\]
Therefore,
\[
\boxed{
p_i=\frac{\partial F_2}{\partial q^i},
\qquad
Q^i=\frac{\partial F_2}{\partial P_i}.
}
\]
This is the type-2 generating function commonly used in mechanics.
Analogously, we can exchange the other variables.
The four conventional types of generating functions

\noindent
\begin{tabular}{cc}
    Function&	Defining relations\\
    \(F_1(q,Q)\)&	\(p=F_{1,q},\quad P=-F_{1,Q}\)\\
    \(F_2(q,P)\)&	\(p=F_{2,q},\quad Q=F_{2,P}\)\\
    \(F_3(p,Q)\)&	\(q=-F_{3,p},\quad P=-F_{3,Q}\)\\
    \(F_4(p,P)\)&	\(q=-F_{4,p},\quad Q=F_{4,P}\)\\
\end{tabular}

The functions are related, wherever the corresponding changes of variables can
be inverted, by
\[
\begin{aligned}
F_2&=F_1+Q^iP_i,\\
F_3&=F_1-q^ip_i,\\
F_4&=F_1-q^ip_i+Q^iP_i.
\end{aligned}
\]
These are partial Legendre transformations.
We arrive at an important distinction:
The canonical transformation is a geometric map of phase space. Its different
generating functions are alternative descriptions of the same map, connected by
Legendre transformations.
The word same is essential here. Transforming \(F_1\) into \(F_2\) does not
ordinarily produce a new canonical transformation.
It only changes how the original transformation is represented.

The four types correspond to different choices of independent variables from
the old and new canonical pairs.
There is no fundamental reason to prefer one type over another. The choice
depends on which projection of the transformation graph is locally invertible.
For example, the identity transformation
\[
Q=q,\qquad P=p
\]
has the simple type-2 generating function
\[
\boxed{F_2(q,P)=q^iP_i.}
\]
But it cannot be represented by a regular type-1 function \(F_1(q,Q)\) in a
neighborhood of the identity: \(q\) and \(Q\) are not independent there,
because \(Q=q\).
This does not mean the identity transformation is singular.
It means only that the chosen generating-function representation is unsuitable.
That distinction will become especially important when discussing
Hamilton–Jacobi caustics.

\subsection{Generating functions and Lagrangian submanifolds}
There is a coordinate-independent interpretation that unifies these
constructions.
Let
\[
M=T^*Q
\]
be phase space, equipped with its symplectic form \(\omega\).
Consider the product space
\[
M\times M
\]
with coordinates
\[
(q,p;Q,P).
\]
Give this product the symplectic form
\[
\boxed{
\Omega=\pi_1^*\omega-\pi_2^*\omega,
}
\]
where \(\pi_1\) and \(\pi_2\) project onto the two factors.
Now consider the graph of a transformation \(\Phi\):
\[
\Gamma_\Phi=
\{(z,\Phi(z)):z\in M\}.
\]
The graph has dimension \(2n\), while the product space has dimension \(4n\).
A direct calculation gives
\[
\left.\Omega\right|_{\Gamma_\Phi}
=\omega-\Phi^*\omega.
\]
Therefore,
\[
\boxed{
\Phi\text{ is canonical}
\iff
\Gamma_\Phi\text{ is Lagrangian in }(M\times M,\Omega).
}
\]
Here Lagrangian means a maximal-dimensional submanifold on which the symplectic
form vanishes.
This is a striking geometric result.
A canonical transformation can be understood not primarily as a change of
coordinates, but as a special submanifold of the product of two phase spaces.
Generating functions are local descriptions of this Lagrangian submanifold.

\paragraph{From a function to a Lagrangian submanifold}
There is a related elementary construction.
Given a smooth function \(S(q)\) on configuration space, form its differential
\[
dS=\frac{\partial S}{\partial q^i}dq^i.
\]
The graph of \(dS\) inside \(T^*Q\) is
\[
\Lambda_S=
\left\{
(q,p):
p_i=\frac{\partial S}{\partial q^i}
\right\}.
\]
On this graph,
\[
\theta|_{\Lambda_S}=dS,
\]
and consequently,
\[
\omega|_{\Lambda_S}=-d^2S=0.
\]
Thus
\[
\boxed{
\text{The graph of }dS\text{ is a Lagrangian submanifold.}
}
\]
This is the geometric foundation of generating functions.
It is also the starting point of the modern geometric treatment of
Hamilton–Jacobi theory.
A generating function does not merely encode derivatives. It specifies a
Lagrangian submanifold through those derivatives.

\subsection{ Physical example: a point transformation versus a general
canonical transformation}
Suppose we change configuration coordinates,
\[
Q^i=Q^i(q).
\]
A natural generating function is
\[
\boxed{F_2(q,P)=Q^i(q)P_i.}
\]
The canonical transformation equations give
\[
p_i=\frac{\partial Q^j}{\partial q^i}P_j,
\]
\[
Q^i=\frac{\partial F_2}{\partial P_i}.
\]
Therefore,
\[
\boxed{
P_j=\frac{\partial q^i}{\partial Q^j}p_i.
}
\]
This is exactly the coordinate transformation law for covectors.
It follows that the canonical extension of an ordinary configuration-coordinate
change is the cotangent lift of that change.
But general canonical transformations are much broader.
For example, using dimensionless canonical variables, consider
\[
\boxed{Q=p,\qquad P=-q.}
\]
It is generated by
\[
F_1(q,Q)=qQ,
\]
because
\[
p=\frac{\partial F_1}{\partial q}=Q,
\qquad
P=-\frac{\partial F_1}{\partial Q}=-q.
\]
This transformation exchanges position and momentum.
It has no interpretation as an ordinary coordinate transformation on
configuration space.
Nevertheless,
\[
dQ\wedge dP=dp\wedge(-dq)=dq\wedge dp,
\]
so it is perfectly canonical.
This demonstrates the physical significance of symplectic geometry:
Position and momentum are not intrinsically distinguished by the symplectic
structure alone. They are distinguished by the choice of configuration-space
projection, Hamiltonian, and other physical structures.
The symplectic structure determines which variables are conjugate, but not
which of them must be called position.

\subsection{Time-dependent generating functions and the transformation of the
Hamiltonian}
We now introduce time dependence.
Consider the Hamiltonian action
\[
S=\int_{t_1}^{t_2}
\left(p_i\dot q^i-H(q,p,t)\right)dt.
\]
Let the canonical transformation be generated by
\[
F_2(q,P,t).
\]
The transformation equations remain
\[
p_i=\frac{\partial F_2}{\partial q^i},
\qquad
Q^i=\frac{\partial F_2}{\partial P_i}.
\]
But now
\[
dF_2
=p_i\,dq^i+Q^i\,dP_i
+\frac{\partial F_2}{\partial t}\,dt.
\]
Using
\[
Q^i\,dP_i=d(Q^iP_i)-P_i\,dQ^i,
\]
we obtain
\[
p_i\,dq^i
=
P_i\,dQ^i
+d(F_2-Q^iP_i)
-\frac{\partial F_2}{\partial t}\,dt.
\]
Hence
\[
p_i\,dq^i-H\,dt
=
P_i\,dQ^i-K\,dt
+d(F_2-Q^iP_i),
\]
where
\[
\boxed{
K(Q,P,t)
=
H(q,p,t)+\frac{\partial F_2}{\partial t},
}
\]
with the right-hand side expressed in the new variables.
This extra time derivative is essential.
For a time-independent canonical transformation,
\[
K=H
\]
as functions evaluated at corresponding phase-space points.
For a time-dependent transformation, the new Hamiltonian differs by \(\partial_tF_2\).
The endpoint term does not affect the local equations of motion when the
transformed action is equipped with the corresponding boundary conditions.

\paragraph{An important physical implication} A Hamiltonian is not invariant
under general time-dependent canonical transformations. Indeed, it is possible
to transform a nonzero Hamiltonian into zero. This does not eliminate the
physical dynamics. It transfers the time dependence from the Hamiltonian into
the transformation connecting the old and new variables.

Example: a free particle transformed into a system with zero Hamiltonian
Consider
\[
H(q,p)=\frac{p^2}{2m}.
\]
Choose the time-dependent type-2 generating function
\[
\boxed{
F_2(q,P,t)=qP-\frac{P^2t}{2m}.
}
\]
Then
\[
p=\frac{\partial F_2}{\partial q}=P,
\]
and
\[
Q=\frac{\partial F_2}{\partial P}
=q-\frac{Pt}{m}.
\]
Thus
\[
\boxed{
P=p,\qquad Q=q-\frac{pt}{m}.
}
\]
The new Hamiltonian is
\[
K
=H+\frac{\partial F_2}{\partial t}
=\frac{P^2}{2m}-\frac{P^2}{2m}.
\]
Therefore,
\[
\boxed{K=0.}
\]
Hamilton's equations in the new variables are simply
\[
\dot Q=0,\qquad \dot P=0.
\]
The transformed variables are constants of motion.
Returning to the original variables,
\[
q(t)=Q+\frac{Pt}{m},
\qquad
p(t)=P.
\]
We have recovered the ordinary free-particle motion.

The transformation has replaced the instantaneous position \(q\) by the
position \(Q\) extrapolated back to \(t=0\), and the instantaneous momentum by
its conserved value. The motion has not disappeared. It has been encoded in the
time-dependent relation between the two sets of variables. This construction
leads directly to Hamilton–Jacobi theory.

\subsection{Hamilton–Jacobi theory as the search for a special canonical transformation}
In Hamilton–Jacobi theory, we try to find a time-dependent canonical
transformation for which the new Hamiltonian vanishes:
\[
K=0.
\]
Let its generating function be
\[
F_2(q,P,t)=S(q,P,t).
\]
Since
\[
K=H(q,p,t)+\frac{\partial S}{\partial t},
\]
and
\[
p_i=\frac{\partial S}{\partial q^i},
\]
the condition \(K=0\) gives
\[
\boxed{
\frac{\partial S}{\partial t}
+
H\left(q,\frac{\partial S}{\partial q},t\right)
=0.
}
\]
This is the Hamilton–Jacobi equation.
The remaining canonical transformation equations are
\[
\boxed{
p_i=\frac{\partial S}{\partial q^i},
\qquad
Q^i=\frac{\partial S}{\partial P_i}.
}
\]
If \(S(q,P,t)\) is a complete integral with
\[
\det
\left(
\frac{\partial^2 S}{\partial q^i\partial P_j}
\right)\ne0,
\]
then the relations can be inverted locally, and \(Q,P\) provide canonical
constants of motion.

The geometric meaning of Hamilton–Jacobi theory is therefore very concrete:
Hamilton–Jacobi theory seeks a generating function for a time-dependent
canonical transformation that converts the dynamics into constant phase-space
coordinates. This explains both its power and its limitations.

\paragraph{The distinction between a complete integral and an arbitrary solution}
Not every solution of the Hamilton–Jacobi equation defines a canonical
transformation of the full phase space.
A particular solution \(S(q,t)\) defines, where smooth, a momentum field
\[
p_i(q,t)=\frac{\partial S}{\partial q^i}.
\]
It therefore specifies a Lagrangian submanifold of phase space, not necessarily
an entire canonical coordinate system.
To obtain a complete canonical transformation, we require a sufficiently rich
family
\[
S(q,P,t)
\]
depending on \(n\) independent parameters \(P_i\) and satisfying the
appropriate nondegeneracy condition.
This distinction is crucial when discussing initial-value problems, multivalued
actions, and caustics.

\subsection{The action as a generating function of time evolution}
There is another extremely important connection.
Consider a regular Lagrangian system and let
\[
S(q,t;Q,t_0)
\]
be the classical action evaluated along a classical trajectory joining the
initial configuration \(Q\) at time \(t_0\) to the final configuration \(q\) at
time \(t\).
That is,
\[
S(q,t;Q,t_0)
=
\int_{t_0}^{t}L(q_{\mathrm{cl}},\dot q_{\mathrm{cl}},s)\,ds.
\]
Assume for the moment that the relevant trajectory exists uniquely and depends
smoothly on the endpoints.
The first variation of the action, evaluated on a classical trajectory, reduces
to boundary terms:
\[
\delta S
=p_i(t)\,\delta q^i
-P_i(t_0)\,\delta Q^i
-H(t)\,\delta t
+H(t_0)\,\delta t_0.
\]
Thus,
\[
\boxed{
\frac{\partial S}{\partial q^i}=p_i(t),
\qquad
\frac{\partial S}{\partial Q^i}=-P_i(t_0).
}
\]
For fixed endpoint times, this has exactly the form of a type-1 generating
function.
The endpoint action therefore generates the canonical relation between initial
and final phase-space data.

Free-particle example
For a free particle moving from \(Q\) at \(t_0=0\) to \(q\) at time \(t>0\),
\[
q_{\mathrm{cl}}(s)=Q+\frac{q-Q}{t}s.
\]
The velocity is constant:
\[
v=\frac{q-Q}{t}.
\]
Therefore,
\[
\boxed{
S(q,t;Q,0)=\frac{m(q-Q)^2}{2t}.
}
\]
Differentiating with respect to the final endpoint gives
\[
p=\frac{\partial S}{\partial q}
=\frac{m(q-Q)}{t}.
\]
Differentiating with respect to the initial endpoint gives
\[
P=-\frac{\partial S}{\partial Q}
=\frac{m(q-Q)}{t}.
\]
Hence
\[
p=P,
\qquad
q=Q+\frac{Pt}{m}.
\]
The usual free-particle trajectory has been reconstructed entirely from
derivatives of its classical action. Notice something remarkable: the action,
initially introduced as an integral along a trajectory, is also a function that
generates the corresponding canonical map. This connects Hamilton's variational
principle directly to symplectic geometry.

\subsection{ Canonical transformations, caustics, and the failure of generating functions}
This connection also clarifies why Hamilton–Jacobi theory encounters caustics
even when Hamiltonian trajectories remain perfectly regular.
Consider the one-dimensional harmonic oscillator,
\[
H(q,p)=\frac{p^2}{2m}+\frac12m\omega^2q^2.
\]
Its phase-space flow is
\[
\begin{pmatrix}q(t)\\p(t)\end{pmatrix}
=
\begin{pmatrix}
\cos\omega t & \dfrac{\sin\omega t}{m\omega}\\
-m\omega\sin\omega t & \cos\omega t
\end{pmatrix}
\begin{pmatrix}Q\\P\end{pmatrix}.
\]
This is a perfectly regular symplectic transformation for all times.
The determinant of the matrix is identically one.
However, for \(\sin\omega t\ne0\), the endpoint action is
\[
\boxed{
S(q,t;Q,0)
=
\frac{m\omega}{2\sin\omega t}
\left[
(q^2+Q^2)\cos\omega t-2qQ
\right].
}
\]
Its mixed second derivative is
\[
\frac{\partial^2 S}{\partial q\,\partial Q}
=-\frac{m\omega}{\sin\omega t}.
\]
At times
\[
t=\frac{k\pi}{\omega},
\]
the endpoint generating function becomes singular.
But the underlying Hamiltonian flow remains nonsingular.
For example, at \(t=\pi/\omega\),
\[
q=-Q,\qquad p=-P.
\]
The canonical transformation still exists perfectly well. What fails is the
attempt to determine the initial momentum uniquely from both endpoint
positions. At this time, the final position is \(-Q\), independent of initial
momentum. This is a failure of a particular projection of the phase-space
transformation, not a failure of Hamiltonian mechanics. In geometric language,
the Lagrangian graph of the canonical transformation is smooth, but its
projection onto the endpoint configuration variables ceases to be locally
invertible. That is the geometric origin of the singularity of this type of
generating function. Relation to caustics For a general Hamilton–Jacobi
solution, caustics occur when the projection of an evolved Lagrangian manifold
onto configuration space ceases to be locally regular. At such points, an
initially single-valued phase function may become multivalued, or its smooth
representation may break down. The full phase-space evolution can still remain
regular.

This gives a useful hierarchy:
\[
\boxed{
\begin{gathered}
\text{Regular Hamiltonian flow}\\
\not\Rightarrow\\
\text{Globally regular endpoint generating function}\\
\not\Rightarrow\\
\text{Globally single-valued Hamilton--Jacobi phase}.
\end{gathered}
}
\]
This distinction is also central to semiclassical mechanics and the emergence
of Maslov indices.

\subsection{ Composition of canonical transformations and the stationary-action
principle}
There is another connection that is not always emphasized in elementary
treatments.
Suppose we have two canonical transformations,
\[
(q,p)\longrightarrow(u,r)
\longrightarrow(Q,P),
\]
generated by type-1 functions
\[
F_a(q,u),\qquad F_b(u,Q).
\]
The first transformation gives
\[
p=\frac{\partial F_a}{\partial q},
\qquad
r=-\frac{\partial F_a}{\partial u}.
\]
The second gives
\[
r=\frac{\partial F_b}{\partial u},
\qquad
P=-\frac{\partial F_b}{\partial Q}.
\]
The intermediate momentum must be the same in both descriptions.
Therefore,
\[
-\frac{\partial F_a}{\partial u}
=
\frac{\partial F_b}{\partial u}.
\]
Or
\[
\boxed{
\frac{\partial}{\partial u}
\left[F_a(q,u)+F_b(u,Q)\right]=0.
}
\]
The generating function of the composed transformation is consequently obtained
by eliminating \(u\) through stationarity:
\[
\boxed{
F(q,Q)=
\operatorname{stat}_{u}
\left[F_a(q,u)+F_b(u,Q)\right].
}
\]
This is a general symplectic composition law.
Now apply it to classical actions.
Suppose a trajectory runs from \(Q\) at time \(t_0\), through \(u\) at an
intermediate time \(t_1\), to \(q\) at time \(t_2\).
Then the endpoint action satisfies, locally and for the relevant classical branches,
\[
\boxed{
S(q,t_2;Q,t_0)
=
\operatorname{stat}_{u}
\left[
S(q,t_2;u,t_1)
+
S(u,t_1;Q,t_0)
\right].
}
\]

The stationarity condition enforces continuity of canonical momentum across the
intermediate point. This is closely related to Hamilton's principle, but here
it emerges directly from the composition of generating functions. It is also
the classical counterpart of the stationary-phase composition of semiclassical
propagators. The geometric meaning is that composition of canonical relations
is represented by elimination of intermediate variables through a
stationary-value operation. Note that the stationary value need not be a
minimum. General classical trajectories correspond to stationary action, not
necessarily least action.

\subsection{Gauge transformations as canonical transformations}
There is a particularly instructive connection between generating functions,
total derivatives, and electromagnetic gauge transformations.
In Gaussian units, the Lagrangian of a charged particle is
\[
L(q,\dot q,t)
=
\frac12m\dot{\mathbf q}^{\,2}
+\frac ec\mathbf A(\mathbf q,t)\cdot\dot{\mathbf q}
-e\phi(\mathbf q,t).
\]
Under the gauge transformation
\[
\mathbf A'=\mathbf A+\nabla\Lambda,
\]
\[
\phi'=\phi-\frac1c\frac{\partial\Lambda}{\partial t},
\]
the Lagrangian becomes
\[
L'=L+\frac ec\frac{d\Lambda}{dt}.
\]
Thus the equations of motion remain unchanged.
But the canonical momentum changes:
\[
\mathbf p=m\dot{\mathbf q}+\frac ec\mathbf A,
\]
\[
\boxed{
\mathbf p'=\mathbf p+\frac ec\nabla\Lambda.
}
\]
This momentum shift is a time-dependent canonical transformation.
Let the old variables be \((q,p)\) and the new ones \((Q,P)\). Choose
\[
\boxed{
F_2(q,P,t)
=q^iP_i-\frac ec\Lambda(q,t).
}
\]
Then
\[
p_i=\frac{\partial F_2}{\partial q^i}
=P_i-\frac ec\partial_i\Lambda,
\]
and
\[
Q^i=\frac{\partial F_2}{\partial P_i}=q^i.
\]
Thus
\[
\boxed{
Q=q,\qquad
P=p+\frac ec\nabla\Lambda.
}
\]
The transformed Hamiltonian is
\[
K=H+\frac{\partial F_2}{\partial t}
=H-\frac ec\frac{\partial\Lambda}{\partial t},
\]
with the old momentum expressed as \(p=P-(e/c)\nabla\Lambda\).
This is exactly the Hamiltonian constructed from the gauge-transformed
potentials.
The combination
\[
\mathbf p-\frac ec\mathbf A
\]
is gauge invariant, even though the canonical momentum itself is not.
This example unites several ideas:
The addition of a total derivative changes the Lagrangian but not the physical
trajectories. It changes the canonical momentum, induces a canonical
transformation of phase space, and modifies the Hamiltonian by a time
derivative of its generating function.
The Legendre transformation and the gauge-induced canonical transformation are
distinct operations, but they fit together consistently.

\subsection{A deeper relationship: contact geometry}
There is also a geometric reason why Legendre transformations and generating
functions look so similar.
For a function \(f(x)\), introduce the contact space with coordinates
\[
(x^i,p_i,z)
\]
and contact one-form
\[
\boxed{\alpha=dz-p_i\,dx^i.}
\]
The graph of a function and its derivatives,
\[
z=f(x),\qquad p_i=\partial_i f,
\]
satisfies
\[
\alpha=0.
\]
Now introduce new coordinates
\[
X_i=p_i,\qquad
\Pi^i=-x^i,\qquad
Z=z-p_ix^i.
\]
Then
\[
dZ-\Pi^i\,dX_i
=
dz-p_i\,dx^i
=\alpha.
\]
Thus the transformation preserves the contact one-form.
If we write
\[
g(p)=p_ix^i-f(x),
\]
then
\[
Z=-g(p).
\]
So the ordinary Legendre transformation can be realized as a contact
transformation that exchanges the roles of variables and slopes.
This is not merely a formal analogy.
In thermodynamics, one uses contact geometry to organize conjugate variables such as
\[
(S,T),\qquad (V,-P).
\]

In mechanics, the corresponding symplectic structures organize canonical
position–momentum pairs. In Hamilton–Jacobi theory, contact geometry naturally
describes first derivatives of phase functions. Thus both Legendre
transformations and canonical generating functions belong to a wider geometric
family involving contact and symplectic structures.

\subsection{ Singularities and global limitations}
One should distinguish three types of problems that elementary treatments
sometimes conflate.
The ordinary Legendre map of a Lagrangian may fail to be invertible because
\[
\det\frac{\partial^2L}{\partial v^i\partial v^j}=0.
\]
This is a singularity of the velocity-to-momentum correspondence and leads
naturally to constrained mechanics.

A particular generating function may fail to describe a canonical
transformation because the projection onto its chosen independent variables is
singular. The harmonic-oscillator example demonstrates that the canonical map
itself may remain smooth.
Finally, a canonical transformation may be globally symplectic without
admitting a globally single-valued generating function of the ordinary kind.
For example, on the cotangent bundle of the circle, with angle \(\varphi\),
consider
\[
\varphi'=\varphi,\qquad p'=p+c.
\]
This preserves
\[
d\varphi\wedge dp.
\]
However,
\[
p\,d\varphi-p'\,d\varphi=-c\,d\varphi
\]
is closed but not globally exact for \(c\ne0\), because its integral around the
circle is nonzero.
Therefore the transformation is symplectic but not globally exact symplectic.
A locally defined generating function exists, but an ordinary globally
single-valued generating function for the difference of canonical one-forms
does not.
This distinction between local and global symplectic geometry becomes important
in mechanics on nontrivial manifolds, semiclassical quantization, and geometric
phases.

\subsection{What is the unifying physical idea?}
We can now connect all these developments.
The passage from a Lagrangian to a Hamiltonian is a Legendre transformation:
\[
L(q,v)\quad\longrightarrow\quad H(q,p).
\]
It replaces velocity variables by momentum variables and, for a regular
Lagrangian, transfers the dynamics to cotangent-bundle geometry.
A canonical transformation acts within Hamiltonian mechanics:
\[
(q,p)\quad\longrightarrow\quad(Q,P).
\]
Its defining property is preservation of the symplectic structure.
A generating function encodes this canonical transformation by derivatives.
Different choices of generating variables are related by partial Legendre
transformations.
Finally, Hamilton–Jacobi theory searches for a special time-dependent canonical
transformation that trivializes the Hamiltonian evolution, at least locally.
These ideas form the sequence
\[
\boxed{
\begin{gathered}
\text{Lagrangian dynamics on }TQ\\
\Big\downarrow\quad\text{Legendre map}\\
\text{Hamiltonian dynamics on }T^*Q\\
\Big\downarrow\quad\text{Canonical transformation}\\
\text{New Hamiltonian variables }(Q,P)\\
\Big\uparrow\quad\text{Generating function}\\
\text{Hamilton--Jacobi equation when }K=0.
\end{gathered}
}
\]
The distinctive feature of Hamiltonian mechanics is that the geometry of phase
space is universal: the symplectic form does not depend on the particular
Hamiltonian.
Generating functions express this geometric structure in terms of scalar
functions, while Legendre transformations allow different descriptions of the
same geometric information.

\section{Why Symplectic Geometry Rather Than Metric Geometry in Hamiltonian
Mechanics?}
The fundamental reason is that metric geometry describes lengths, angles, and
distances, whereas symplectic geometry describes the pairing of dynamical
variables and the structure of conservative evolution. A Hamiltonian does not
ordinarily generate motion in the direction in which energy increases most
rapidly. Instead, for an autonomous system, it generates motion along surfaces
of constant energy. A metric naturally converts the differential of a function
into its gradient, which is perpendicular to its level surfaces. A symplectic
form converts the same differential into a Hamiltonian vector field, which is
tangent to those surfaces. This difference is not merely mathematical. It
distinguishes conservative Hamiltonian dynamics from gradient dynamics, such as
dissipative relaxation.

\paragraph{But there is a deeper question:}
Why does the variational formulation of mechanics lead naturally to an
antisymmetric two-form rather than a symmetric metric tensor? The answer
connects the order of the equations of motion, the structure of the action, the
Legendre transformation, canonical transformations, and the geometry of phase
space.

\subsection{The distinction between a gradient flow and a Hamiltonian flow}
Let \(M\) be a smooth manifold representing the states of a physical system,
and let
\[
H:M\longrightarrow\mathbb R
\]
be a smooth function interpreted as energy.
The differential \(dH\) is a covector. To obtain equations of motion, we need a
rule that converts this covector into a vector field:
\[
dH\longmapsto X_H.
\]
A metric and a symplectic form provide two fundamentally different rules.
Metric geometry
Given a positive-definite Riemannian metric \(g\), the gradient of \(H\) is
defined by
\[
\boxed{g(\operatorname{grad}_g H,Y)=dH(Y)}
\]
for every vector field \(Y\).
If we choose the dynamics
\[
\dot z=-\operatorname{grad}_gH,
\]
then
\[
\begin{aligned}
\frac{dH}{dt}
&=dH(\dot z)\\
&=-g(\operatorname{grad}_gH,\operatorname{grad}_gH).
\end{aligned}
\]
Therefore,
\[
\boxed{\frac{dH}{dt}\leq0.}
\]
The motion proceeds downhill in energy.
This is precisely the structure expected for a relaxation process.
Symplectic geometry
Instead, suppose \(M\) is equipped with a nondegenerate antisymmetric two-form
\(\omega\).
The Hamiltonian vector field is defined by
\[
\boxed{\iota_{X_H}\omega=dH.}
\]
In other words,
\[
\omega(X_H,Y)=dH(Y).
\]
Taking \(Y=X_H\),
\[
dH(X_H)=\omega(X_H,X_H)=0,
\]
because \(\omega\) is antisymmetric.
Consequently,
\[
\boxed{\frac{dH}{dt}=0.}
\]
The Hamiltonian flow preserves the energy without requiring a metric.
That is the first essential distinction. 
Of course, energy conservation alone does not uniquely determine Hamiltonian
motion. There are many vector fields tangent to energy surfaces. The symplectic
structure supplies the additional information needed to select a particular
one.

\subsection{ Why an antisymmetric tensor appears}
We can derive the need for antisymmetry from fairly general physical
assumptions.
Suppose we want the equations of motion to be first order,
\[
\dot z^a=X^a(z),
\]
and suppose that the vector field depends linearly on the differential of the
Hamiltonian:
\[
\boxed{\dot z^a=B^{ab}(z)\partial_bH.}
\]
Here \(B^{ab}\) is a contravariant two-tensor.
The rate of change of energy is
\[
\frac{dH}{dt}
=\partial_aH\,B^{ab}\partial_bH.
\]
If we require energy conservation for every possible Hamiltonian, then
\[
\partial_aH\,B^{ab}\partial_bH=0
\]
for every covector \(dH\).
This implies that the symmetric part of \(B\) vanishes:
\[
\boxed{B^{ab}=-B^{ba}.}
\]
Thus an antisymmetric structure is the natural candidate for conservative
dynamics under these assumptions.
If \(B\) is nondegenerate, it has an inverse that defines an antisymmetric
two-form \(\omega\), with the appropriate index and sign conventions.
The relation between energy and velocity becomes
\[
\boxed{\iota_{X_H}\omega=dH.}
\]
We have therefore recovered the fundamental Hamiltonian relation.

\paragraph{Why must the symplectic form also be closed?}
Antisymmetry ensures energy conservation, but it does not yet give us the full
symplectic structure.
A symplectic form must satisfy
\[
\boxed{d\omega=0.}
\]
This property has two closely related meanings.
First, it ensures that Hamiltonian evolution preserves \(\omega\). Second, it
guarantees that the associated Poisson bracket satisfies the Jacobi identity.
Indeed, by Cartan's identity,
\[
\mathcal L_{X_H}\omega
=d(\iota_{X_H}\omega)+\iota_{X_H}d\omega.
\]
Using
\[
\iota_{X_H}\omega=dH,\qquad d\omega=0,
\]
we obtain
\[
\boxed{
\mathcal L_{X_H}\omega=d^2H=0.
}
\]
Consequently, if \(\Phi_t\) is the Hamiltonian flow,
\[
\boxed{\Phi_t^*\omega=\omega.}
\]
This is a much stronger condition than energy conservation.It says that the
fundamental geometric structure relating positions and momenta remains
unchanged under evolution. There is a qualification: this argument derives the
antisymmetric character from energy conservation, but it does not prove that
all conservative systems must be symplectic. Noncanonical Poisson systems,
including degenerate ones, also occur. What distinguishes regular Hamiltonian
mechanics is the combination of nondegeneracy, closure, and the variational
structure.
\subsection{ Why the phase-space symplectic form is canonical}
We now arrive at what I regard as the most convincing argument.
Let \(Q\) be the configuration manifold of a mechanical system.
Coordinates and velocities belong to the tangent bundle,
\[
(q,v)\in TQ,
\]
whereas coordinates and momenta belong to the cotangent bundle,
\[
(q,p)\in T^*Q.
\]
There is a natural pairing between covectors and vectors,
\[
\langle p,v\rangle=p_iv^i.
\]
This pairing exists without any metric.
In particular, \(p_i\) and \(v^i\) are not related by index raising or lowering
unless an additional structure is supplied.
The canonical one-form
On \(T^*Q\), there is a geometrically defined one-form
\[
\boxed{\theta=p_i\,dq^i.}
\]
To define it without coordinates, let
\[
\pi:T^*Q\to Q
\]
be the natural projection.
At a point \((q,p)\in T^*Q\), and for a tangent vector \(W\in T_{(q,p)}T^*Q\), set
\[
\boxed{\theta_{(q,p)}(W)=p(T\pi(W)).}
\]
That is, project the phase-space displacement onto a configuration-space
displacement and let the momentum covector act on it.
Nothing else is required.
No mass, no potential, no connection, and no metric.
Now define
\[
\boxed{\omega=-d\theta.}
\]
In canonical coordinates,
\[
\boxed{\omega=dq^i\wedge dp_i.}
\]
The exterior derivative automatically gives
\[
d\omega=0.
\]
Nondegeneracy follows immediately from the canonical coordinate expression.

This construction works on the cotangent bundle of any smooth configuration
manifold, even when the manifold is curved, has nontrivial topology, or has no
globally defined Euclidean coordinates. Thus the cotangent bundle possesses a
symplectic structure before any particular mechanical system has been
specified. That is not true of a metric on phase space. Such a metric requires
further choices.

\subsection{The symplectic form comes directly from the action principle}
There is an even more physical way to see its origin.
Take the first-order Hamiltonian action
\[
\boxed{
S[q,p]=\int_{t_1}^{t_2}
\left(p_i\dot q^i-H(q,p,t)\right)dt.
}
\]
Varying independently \(q^i\) and \(p_i\), we obtain
\[
\begin{aligned}
\delta S
={}&[p_i\delta q^i]_{t_1}^{t_2}\\
&+\int_{t_1}^{t_2}
\left[
(\dot q^i-H_{p_i})\delta p_i
-(\dot p_i+H_{q^i})\delta q^i
\right]dt.
\end{aligned}
\]
The bulk terms give Hamilton's equations.
But notice the boundary term,
\[
\boxed{p_i\delta q^i.}
\]
This is the evaluation of the canonical one-form \(\theta\) on an endpoint
variation.
Its exterior derivative is precisely what produces the symplectic structure.
The phase-space action therefore naturally contains two geometrically distinct
ingredients:
\[
\boxed{
S=\int\theta-\int H\,dt.
}
\]
The first term is determined by the canonical pairing of position and momentum.
The second contains the particular dynamics.
This separation is extremely important.
The symplectic structure tells us how a Hamiltonian generates a flow. The
Hamiltonian tells us which flow occurs.
It is analogous, though not identical, to the distinction between a geometric
structure and a physical field defined on that structure.

\paragraph{Deriving symplectic geometry from a regular Lagrangian}
Start instead from an autonomous Lagrangian,
\[
L(q^i,v^i).
\]
Its Legendre map is
\[
p_i=\frac{\partial L}{\partial v^i}.
\]
Pulling back the canonical one-form gives the Lagrangian one-form
\[
\boxed{
\theta_L=
\frac{\partial L}{\partial v^i}dq^i.
}
\]
Define
\[
\boxed{\omega_L=-d\theta_L.}
\]
The Lagrangian energy is
\[
\boxed{
E_L=v^i\frac{\partial L}{\partial v^i}-L.
}
\]
If the velocity Hessian is nonsingular, \(\omega_L\) is symplectic.
Let \(\Gamma\) be the second-order dynamical vector field on \(TQ\),
\[
\Gamma=v^i\frac{\partial}{\partial q^i}
+a^i(q,v)\frac{\partial}{\partial v^i}.
\]
Then the Euler–Lagrange equations are equivalent to
\[
\boxed{\iota_\Gamma\omega_L=dE_L.}
\]
To verify this, write \(p_i=L_{v^i}\). Then
\[
\omega_L=dq^i\wedge dp_i,
\]
so
\[
\iota_\Gamma\omega_L
=v^i\,dp_i-\Gamma(p_i)\,dq^i.
\]
On the other hand,
\[
dE_L=v^i\,dp_i-L_{q^i}\,dq^i.
\]
Therefore,
\[
\iota_\Gamma\omega_L=dE_L
\]
holds exactly when
\[
\Gamma(p_i)=L_{q^i},
\]
which is the Euler–Lagrange equation.

We have now obtained the symplectic geometry directly from the variational
problem. This is not a structure arbitrarily attached to Newtonian mechanics.
For regular Lagrangian systems, it is intrinsic to the action.

\subsection{ The metric and the symplectic form play different physical roles}
One might now ask why a metric is so prominent in the Lagrangian formulation.
Consider motion on a Riemannian configuration manifold \((Q,g)\) with
\[
L(q,v)=\frac m2g_{ij}(q)v^iv^j-V(q).
\]
The metric determines the kinetic energy:
\[
T=\frac m2g_{ij}v^iv^j.
\]
The Legendre transformation gives
\[
p_i=mg_{ij}v^j,
\]
and the Hamiltonian becomes
\[
\boxed{
H(q,p)=\frac{1}{2m}g^{ij}(q)p_ip_j+V(q).
}
\]
But the symplectic form remains
\[
\boxed{\omega=dq^i\wedge dp_i.}
\]
Hamilton's equations are
\[
\dot q^i=\frac1m g^{ij}p_j,
\]
\[
\dot p_i
=-\frac{1}{2m}\frac{\partial g^{jk}}{\partial q^i}p_jp_k
-\frac{\partial V}{\partial q^i}.
\]
Eliminating momenta yields
\[
\boxed{
\ddot q^i+\Gamma^i{}_{jk}\dot q^j\dot q^k
=-\frac1m g^{ij}\partial_jV.
}
\]
Thus the configuration-space metric determines inertia, geodesics, and the
relation between velocities and momenta.
The symplectic form determines the canonical structure of evolution.
They answer different questions.
A particularly instructive example: geodesic motion
For \(V=0\), the Hamiltonian is
\[
H=\frac1{2m}g^{ij}p_ip_j.
\]
Its Hamiltonian flow on \(T^*Q\) projects onto geodesics of the metric \(g\) on
\(Q\).
So a single physical problem involves both geometries:
\begin{itemize}
\item The Riemannian metric determines which paths are geodesics.
\item The symplectic structure determines how the corresponding Hamiltonian
    flow is constructed in phase space.
\end{itemize}
Changing the metric changes the geodesic dynamics, but it does not change the
canonical symplectic form on \(T^*Q\). The analogous distinction holds in
relativity. A spacetime metric \(g_{\mu\nu}\) determines the mass-shell
relation and geodesic motion, whereas the cotangent bundle of spacetime still
has its canonical symplectic form \(dx^\mu\wedge dp_\mu\).

\subsection{ Why symplectic geometry is more flexible than metric geometry}
Consider canonical coordinates \((q,p)\) in a two-dimensional phase space.
The transformation
\[
\boxed{
Q=\lambda q,\qquad P=\frac p\lambda
}
\]
preserves the symplectic form:
\[
dQ\wedge dP
=\lambda\,dq\wedge\frac1\lambda dp
=dq\wedge dp.
\]
But it does not preserve the Euclidean metric
\[
ds^2=dq^2+dp^2.
\]
Indeed,
\[
dQ^2+dP^2
=\lambda^2dq^2+\frac1{\lambda^2}dp^2.
\]
Unless \(|\lambda|=1\), lengths and angles change.

The transformation can stretch a phase-space region strongly in one direction
and compress it in the conjugate direction. Nevertheless, it remains canonical.
This is physically appropriate because canonical transformations are not
required to preserve distances in phase space. In fact, a Euclidean distance
such as \[ dq^2+dp^2 \] is not even dimensionally meaningful unless one
introduces scales that allow position and momentum to be compared. By contrast,
\[ dq\wedge dp \] has the physical dimensions of action. The symplectic form
describes an oriented area in each conjugate position–momentum plane, not a
notion of distance. This leads to Liouville's theorem.

\subsection{ Liouville's theorem: conservation of phase-space volume}
For a \(2n\)-dimensional symplectic manifold, the symplectic form determines a
natural volume form,
\[
\boxed{
\Omega=\frac{\omega^n}{n!}.
}
\]
Since Hamiltonian evolution preserves \(\omega\), it also preserves \(\Omega\):
\[
\boxed{\Phi_t^*\Omega=\Omega.}
\]
In canonical coordinates, this is the familiar invariance of phase-space volume,
\[
dq^1\cdots dq^n\,dp_1\cdots dp_n.
\]
For the canonical Hamilton equations,
\[
\dot q^i=H_{p_i},\qquad
\dot p_i=-H_{q^i},
\]
the phase-space divergence is
\[
\begin{aligned}
\operatorname{div}X_H
&=\sum_i\left(
\frac{\partial\dot q^i}{\partial q^i}
+\frac{\partial\dot p_i}{\partial p_i}
\right)\\
&=\sum_i\left(
\frac{\partial^2H}{\partial q^i\partial p_i}
-\frac{\partial^2H}{\partial p_i\partial q^i}
\right)\\
&=0.
\end{aligned}
\]
Therefore Hamiltonian flows preserve volume.

But symplectic geometry says more than this: a Hamiltonian flow preserves the
full two-form, not just its highest exterior power. In particular, it preserves
the symplectic pairing of transported tangent vectors. This is why phase-space
volume conservation is a consequence of symplectic structure, rather than its
defining property.

\paragraph{Contrast with friction}
Consider the damped harmonic oscillator in ordinary mechanical variables,
\[
\dot q=\frac pm,\qquad
\dot p=-kq-\gamma p.
\]
Its mechanical energy is
\[
H=\frac{p^2}{2m}+\frac{kq^2}{2}.
\]
Then
\[
\frac{dH}{dt}=-\frac{\gamma p^2}{m}.
\]
Its phase-space divergence is
\[
\operatorname{div}X=-\gamma.
\]
Thus mechanical phase-space areas contract.
For the standard two-form \(\omega=dq\wedge dp\),
\[
\boxed{\mathcal L_X\omega=-\gamma\omega.}
\]

This is not Hamiltonian evolution relative to that fixed symplectic form. The
example makes the physical distinction concrete. Dissipation naturally
contracts phase-space volumes; standard autonomous Hamiltonian dynamics cannot
do so. A damped oscillator can, of course, be embedded in a larger Hamiltonian
system or described using a time-dependent Hamiltonian and suitably modified
momenta. Those constructions change the geometric description and do not
contradict this result.

\subsection{The Poisson bracket: another reason a symmetric metric is
insufficient}
Symplectic geometry produces an antisymmetric bracket between observables.
In canonical coordinates,
\[
\boxed{
\{f,g\}
=
\frac{\partial f}{\partial q^i}
\frac{\partial g}{\partial p_i}
-
\frac{\partial f}{\partial p_i}
\frac{\partial g}{\partial q^i}.
}
\]
The time evolution of any observable \(f\) is
\[
\boxed{
\frac{df}{dt}=\{f,H\}+\frac{\partial f}{\partial t}.
}
\]
The Poisson bracket has three essential properties:
\[
\{f,g\}=-\{g,f\},
\]
\[
\{f,gh\}=g\{f,h\}+h\{f,g\},
\]
and the Jacobi identity,
\[
\boxed{
\{f,\{g,h\}\}
+\{g,\{h,f\}\}
+\{h,\{f,g\}\}=0.
}
\]
The Jacobi identity is particularly important: it ensures the algebraic
consistency of infinitesimal canonical transformations.
There is no analogous natural Poisson bracket obtained merely from a Riemannian
metric.
One can construct a symmetric bracket using a metric,
\[
(f,g)_g=g(\operatorname{grad}f,\operatorname{grad}g),
\]
but it generates a different type of dynamics.
In particular,
\[
(H,H)_g\geq0,
\]
whereas
\[
\{H,H\}=0.
\]
The distinction between symmetric and antisymmetric brackets corresponds
naturally to the distinction between dissipative and conservative processes.
This becomes especially useful in formulations that combine Hamiltonian motion
and dissipation, such as metriplectic dynamics.

\subsection{ Darboux's theorem versus Riemannian curvature}
There is a remarkable difference between metric and symplectic geometry that
has a direct physical interpretation.

In Riemannian geometry, one can always make the metric Euclidean at a point by
choosing appropriate local coordinates, but generally not throughout a
neighborhood.
The obstruction is curvature.
The Riemann tensor contains genuine local geometric information.

Symplectic geometry is different.
Darboux's theorem states that every symplectic manifold is locally
symplectomorphic to the standard canonical model.
Around every point, one can choose local coordinates such that
\[
\boxed{
\omega=\sum_{i=1}^n dq^i\wedge dp_i.
}
\]
There is no local analogue of the Riemann curvature tensor that distinguishes
symplectic forms.
All symplectic manifolds of a fixed dimension are locally equivalent as
symplectic manifolds.

This has an important consequence for physics. If the symplectic form can
always be made canonical locally, then the distinctive local dynamics must
reside in the Hamiltonian, rather than in local curvature invariants of the
symplectic form. Of course, symplectic geometry still has significant global
invariants and topology. Darboux's theorem is only local. 

\subsection{ Can metric and symplectic geometry coexist?}
Yes, and this is a very interesting point.
There is no contradiction between the two structures.
In fact, symplectic manifolds admit compatible almost complex structures and
Riemannian metrics.
Consider, for simplicity, phase space \(\mathbb R^{2n}\) with
\[
\omega=\sum_i dq^i\wedge dp_i.
\]
Choose the Euclidean metric
\[
g=\sum_i(dq^i)^2+(dp_i)^2.
\]
Define an almost complex structure \(J\) by
\[
J\frac{\partial}{\partial q^i}
=\frac{\partial}{\partial p_i},
\qquad
J\frac{\partial}{\partial p_i}
=-\frac{\partial}{\partial q^i}.
\]
Then
\[
J^2=-I.
\]
With these conventions,
\[
\boxed{g(X,Y)=\omega(X,JY).}
\]
The Hamiltonian vector field satisfies
\[
\boxed{
X_H=-J\operatorname{grad}_gH.
}
\]
Thus, in these compatible coordinates, the Hamiltonian vector field is the
metric gradient rotated by a quarter-turn in each canonical plane.
For one degree of freedom,
\[
X_H=
\begin{pmatrix}
0&1\\
-1&0
\end{pmatrix}
\begin{pmatrix}
H_q\\
H_p
\end{pmatrix}.
\]
The metric gradient points toward increasing energy, and the Hamiltonian vector
field is tangent to the energy contours.
But the important point is that \(J\) is additional structure. The symplectic
form alone does not select a unique compatible metric.

This is the beginning of almost Kähler geometry, with Kähler geometry arising
when further integrability conditions hold.
It is also one of the geometric bridges toward quantum mechanics, where
symplectic, metric, and complex structures occur together in the geometry of
Hilbert space.

\subsection{ What happens in systems with magnetic fields?}
A useful test of the preceding interpretation is a charged particle in a
magnetic field.
In Gaussian units,
\[
L(q,v)=
\frac m2g_{ij}(q)v^iv^j
+\frac ecA_i(q)v^i-V(q).
\]
The canonical momentum is
\[
p_i=mg_{ij}v^j+\frac ecA_i.
\]
Define the kinetic momentum
\[
\pi_i=p_i-\frac ecA_i.
\]
In the canonical variables \((q,p)\),
\[
\omega=dq^i\wedge dp_i.
\]
But expressing the same two-form in the variables \((q,\pi)\) gives
\[
\boxed{
\omega
=dq^i\wedge d\pi_i
-\frac e{2c}F_{ij}\,dq^i\wedge dq^j,
}
\]
where
\[
F_{ij}=\partial_iA_j-\partial_jA_i.
\]
The magnetic field now appears explicitly in the symplectic form.
Correspondingly,
\[
\boxed{
\{\pi_i,\pi_j\}=\frac ecF_{ij}.
}
\]
This is especially instructive.
In canonical momentum variables, the symplectic form has its universal
expression and the vector potential appears in the Hamiltonian.
In kinetic momentum variables, the Hamiltonian is simpler, but the magnetic
field appears in the symplectic form and Poisson brackets.
The geometry and Hamiltonian can therefore exchange some of their apparent
complexity under a change of variables.
The magnetic example also demonstrates that the relevant symplectic geometry is
not simply a mathematical replacement for metric geometry. It can encode
interactions through noncanonical variables.

\subsection{Is symplectic geometry inevitable?}
I would distinguish three statements.
\begin{enumerate}
   \item Newton's equations do not explicitly require symplectic geometry in
       order to be written down. One can formulate the dynamics as second-order
        differential equations on configuration space.
   \item A regular Lagrangian variational formulation naturally induces
       symplectic geometry through its Legendre map and canonical one-form.
   \item Once we formulate conservative dynamics in terms of a Hamiltonian on
       phase space, symplectic geometry is the appropriate intrinsic structure
        for expressing the canonical pairing, equations of motion, and their
        invariance.
\end{enumerate}

So the correct statement is not that metric geometry is replaced or superseded
by symplectic geometry. Rather, these geometries serve different functions, and
the variational passage from configuration-space mechanics to phase-space
mechanics naturally produces both.

An important caveat concerns constrained systems. For singular Lagrangians, the
initial geometric structure may be presymplectic (closed but degenerate), with
symplectic geometry emerging only after suitable reduction. More general
Hamiltonian systems can also be described on Poisson manifolds whose brackets
are degenerate.

\subsection{ The broader physical interpretation}
We can now put the relationships in a single table.

\noindent
\begin{tabular}{p{0.3\textwidth}p{0.3\textwidth}p{0.3\textwidth}}
    Geometric structure&	Mathematical object&	Physical role\\
    Riemannian geometry&	Symmetric tensor \(g\)&	Distances, kinetic energy, inertia, geodesics\\
    Symplectic geometry&	Closed antisymmetric form \(\omega\)&	Canonical pairing, conservative Hamiltonian flow\\
    Poisson geometry&	Antisymmetric bivector satisfying Jacobi identity&	Dynamics, including degenerate brackets and Casimirs\\
    Contact geometry&	Contact distribution, locally \(\ker\alpha\)&	Thermodynamic state geometry and certain dissipative formulations\\
    Complex/Kähler geometry&	Compatible \(g,\omega,J\)&	Unified metric, symplectic, and complex structures\\
\end{tabular}

There is also an interesting hierarchy of invariance groups.
Riemannian geometry is characterized by transformations preserving the metric.
Symplectic geometry is characterized by transformations preserving the
symplectic form.
The symplectic group is not the orthogonal group. It admits transformations
that arbitrarily stretch positions while compressing their conjugate momenta.
From the physical point of view, this is natural: a canonical transformation
need not preserve the numerical scale of a coordinate or momentum. It must
preserve their conjugate structure.

\subsection{ Final perspective: what is fundamental?}
The central distinction can be expressed in two equations:
\[
\boxed{
\begin{aligned}
g(\operatorname{grad}_g H,\cdot)&=dH,\\
\omega(X_H,\cdot)&=dH.
\end{aligned}
}
\]
Both structures convert an energy differential into a vector field.
But the first uses a symmetric positive-definite pairing. The second uses an
antisymmetric nondegenerate pairing.
As a consequence,
\[
\boxed{
\begin{aligned}
dH(\operatorname{grad}_gH)&=
g(\operatorname{grad}_gH,\operatorname{grad}_gH),\\
dH(X_H)&=\omega(X_H,X_H)=0.
\end{aligned}
}
\]
The first relation describes the direction of energy increase; the second
describes evolution that conserves energy.
But perhaps the deepest argument is not energy conservation itself. It is the
geometric structure of the variational principle:
\[
\boxed{
\begin{gathered}
\text{Action functional}\\
\Downarrow\\
\text{Boundary term }p_i\,\delta q^i\\
\Downarrow\\
\text{Canonical one-form }\theta=p_i\,dq^i\\
\Downarrow\\
\text{Symplectic form }\omega=-d\theta\\
\Downarrow\\
\text{Hamiltonian dynamics }\iota_{X_H}\omega=dH.
\end{gathered}
}
\]
That derivation explains not merely why symplectic geometry is convenient, but
why it arises intrinsically from regular variational mechanics.

